\documentclass[10pt,a4paper]{amsart}
\usepackage[margin=19mm]{geometry}
\usepackage{amsmath,amssymb,mathtools}
\usepackage[scr=rsfs,cal=euler]{mathalfa}
\usepackage{xfrac}
\usepackage[unicode,hidelinks,bookmarks=false]{hyperref}
\newcommand{\ie}{i.e.\ }
\newcommand{\cf}{cf.\ }
\newcommand{\resp}{respectively\ }
\newcommand{\Subsection}{\subsection}
\providecommand{\nobreakdash}{\nobreak\hbox{-}}
\setlength{\emergencystretch}{2em}
\multlinegap0pt
\usepackage{amssymb}
\usepackage{mathtools}
\usepackage{xcolor}
\usepackage{tikz}
\usepackage[shortlabels]{enumitem}
\newtheorem{lemma}{Lemma}
\newtheorem{proposition}{Proposition}
\newcommand{\E}{\mathbb{E}}
\newcommand{\KL}{\operatorname{KL}}
\newcommand{\thetalo}{\theta_-}
\title{On Observation Time for Recovering \\ Latent Hawkes Networks}
\author{Jonas Linkerh\"agner, Michele Bortolasi, Lorenzo Baldassari, Maarten V. de Hoop, Ivan Dokmani\'c}
\date{Source: arXiv:2605.08400v1 (CC BY 4.0)}
\begin{document}
\maketitle

\noindent\emph{Source and licence.} On Observation Time for Recovering \\ Latent Hawkes Networks. Version arXiv:2605.08400v1 (CC BY 4.0), distributed by arXiv under the Creative Commons Attribution 4.0 International licence, http://creativecommons.org/licenses/by/4.0/, as recorded in the arXiv licence metadata for this exact version and shown on the abstract page https://arxiv.org/abs/2605.08400v1. Full licence notice: \texttt{licenses/arxiv-2605.08400-CC-BY-4.0.txt}.\par\smallskip

\noindent\textit{Editorial scope.} Counted unit: the article's proposition prop:app-initKL-k (source0.tex lines 2316--2329). The article's complete original proof of it is delivered with it (source0.tex lines 2331--2435, one \texttt{proof} environment of 105 lines). No mathematics is written, repaired or re-derived: the statement and the proof are the article's own lines, in the article's own notation, and no retained line is substituted or re-flowed.  Retained uncounted as the source-local context the proof needs: lines 2263--2278.  Nothing was omitted from any retained span, so the package carries no omission record.  No retained line contains a citation, so the wrapper carries no bibliography and no bibliography entry.  No retained line contains a figure environment, an image-inclusion command or drawing code, so no image is delivered and the package declares no asset dependency.  Editorial formatting only, in addition: an AMS \texttt{amsart} preamble that restates the article's own declarations and macros used by the retained lines, namely the environments \texttt{lemma}, \texttt{proposition} on separate counters, together with the standard packages those lines need.  The retained environments print in this document as Lemma 1, Proposition 1 in order of appearance, whereas the article prints the same items with the numbering of its own full document.  The complete native source of the article as distributed in the arXiv e-print for this exact version is bundled unchanged as \texttt{sources/200151/source0.tex}.
\begin{lemma}\label{lem:app-translation-kl}
Let $q_0:=q_{\bar\mu_\star}$ be the density of $\pi_{\bar\mu_\star}^{\mathrm P}$.
For $w\ge 0$, define the translated density
\[
q_{0,w}(z):=q_0(z-w)\,\mathbf 1_{\{z>w\}},
\qquad z>0,
\]
that is, the law of $U+w$ when $U\sim \pi_{\bar\mu_\star}^{\mathrm P}$.
Then there exists a finite constant $C_{\mathrm{tr}}$, depending only on
$(\beta,\bar\mu_\star)$, such that
\[
g(w):=\KL(q_{0,w}\,dz\|q_0\,dz)
\le C_{\mathrm{tr}}\bigl(1+w\log(1+w)\bigr),
\qquad w\ge 0.
\]
\end{lemma}

\begin{proposition}
\label{prop:app-initKL-k}
For every $d\ge k+2$ and every $S\in\mathcal S_{d,k}(i_\star)$,
\[
\KL(\pi_S\|\pi_0)
=
\KL\!\bigl(\bar\pi_k\|\pi_{\bar\mu_\star}^{\mathrm P}\otimes
(\pi_{\bar\mu}^{\mathrm P})^{\otimes k}\bigr)
=:C_{\mathrm{init}}.
\]
The constant $C_{\mathrm{init}}$ is finite and depends only on
$(\beta,\bar\mu,\bar\mu_\star,\thetalo,k)$; in particular it is independent of
both $S$ and $d$.
\end{proposition}

\begin{proof}
Fix $S=\{s_1,\dots,s_k\}\in\mathcal S_{d,k}(i_\star)$ and order the remaining
coordinates in any fixed way. Let $R_S$ be the coordinate permutation that puts
$i_\star$ first, then the elements of $S$, and then all other coordinates. KL is
invariant under this bijection.

Under $\Theta^{(S)}$, all coordinates outside $\{i_\star\}\cup S$ are independent
Poisson shot-noise coordinates with stationary law $\pi_{\bar\mu}^{\mathrm P}$,
and they are independent of the target/source block. The nontrivial block
$(Z_t,\mathbf Y_t):=(X_{i_\star}(t),(X_{s_1}(t),\dots,X_{s_k}(t)))$ satisfies
\[
\begin{aligned}
dY_t^{(m)}&=-\beta Y_t^{(m)}\,dt+dN_t^{(m)},\qquad m=1,\dots,k,\\
dZ_t&=-\beta Z_t\,dt+dM_t,
\end{aligned}
\]
where $N^{(1)},\dots,N^{(k)}$ are independent Poisson processes of rate
$\bar\mu$ and $M$ has predictable intensity
$\bar\mu_\star+\thetalo\sum_{m=1}^kY_{t-}^{(m)}$. If $\bar\pi_k$ denotes the
stationary law of this block, then
\[
R_S\#\pi_S=\bar\pi_k\otimes(\pi_{\bar\mu}^{\mathrm P})^{\otimes(d-k-1)}.
\]
Under the baseline model,
\[
R_S\#\pi_0=\pi_{\bar\mu_\star}^{\mathrm P}\otimes
(\pi_{\bar\mu}^{\mathrm P})^{\otimes k}\otimes
(\pi_{\bar\mu}^{\mathrm P})^{\otimes(d-k-1)}.
\]
Product additivity of relative entropy gives
\[
\KL(\pi_S\|\pi_0)=
\KL\!\bigl(\bar\pi_k\|\pi_{\bar\mu_\star}^{\mathrm P}\otimes
(\pi_{\bar\mu}^{\mathrm P})^{\otimes k}\bigr),
\]
and the right-hand side is independent of the particular set $S$ by
exchangeability of the source coordinates.

It remains only to prove finiteness. Let
$\mathcal H=\sigma(Y_s^{(m)}:s\le0,1\le m\le k)$ be the stationary source
history. Conditional on $\mathcal H$, the target process on $(-\infty,0]$ is an
inhomogeneous Poisson process with intensity
$\bar\mu_\star+\thetalo\sum_mY_{t-}^{(m)}$. Decompose it as $M=M^{(0)}+M^{(1)}$,
where, conditional on $\mathcal H$, $M^{(0)}$ is homogeneous Poisson of rate
$\bar\mu_\star$, $M^{(1)}$ is Poisson with intensity
$\thetalo\sum_mY_{t-}^{(m)}\,dt$, and the two are independent. Then
\[
Z_0=U+W,
\qquad
U:=\int_{(-\infty,0]}e^{\beta s}\,dM_s^{(0)},\qquad
W:=\int_{(-\infty,0]}e^{\beta s}\,dM_s^{(1)}.
\]
Here $U\sim\pi_{\bar\mu_\star}^{\mathrm P}$ and $U$ is independent of
$(\mathcal H,W)$.

The $\mathbf Y_0$-marginal of $\bar\pi_k$ is
$(\pi_{\bar\mu}^{\mathrm P})^{\otimes k}$. Hence the entropy chain rule and
conditional Jensen's inequality give
\[
\KL\!\bigl(\bar\pi_k\|\pi_{\bar\mu_\star}^{\mathrm P}\otimes
(\pi_{\bar\mu}^{\mathrm P})^{\otimes k}\bigr)
= \E\bigl[\KL(P_{Z_0\mid \mathbf Y_0}\|\pi_{\bar\mu_\star}^{\mathrm P})\bigr]
\le \E\bigl[\KL(P_{Z_0\mid \mathcal H}\|\pi_{\bar\mu_\star}^{\mathrm P})\bigr].
\]
Indeed, since \(\mathbf Y_0\) is \(\mathcal H\)-measurable,
\(P_{Z_0\mid \mathbf Y_0}\) is obtained by averaging
\(P_{Z_0\mid \mathcal H}\) over the conditional law of \(\mathcal H\) given
\(\mathbf Y_0\), and relative entropy is convex in its first argument.
Given $\mathcal H$, the conditional law of $Z_0$ is the mixture of translates
$q_{0,w}(z)\,dz$ with mixing law $P_{W\mid\mathcal H}$. Convexity of KL and
Lemma~\ref{lem:app-translation-kl} therefore yield
\[
\KL(P_{Z_0\mid \mathcal H}\|\pi_{\bar\mu_\star}^{\mathrm P})
\le \E[g(W)\mid\mathcal H]
\le C\, \E[1+W\log(1+W)\mid\mathcal H].
\]
Thus finiteness follows once $\E[W\log(1+W)]<\infty$.

In fact $W$ has exponential moments of all orders. For $\eta>0$, the conditional
Laplace functional gives
\[
\E[e^{\eta W}\mid\mathcal H]
=
\exp\!\left(\thetalo\int_{-\infty}^0
(e^{\eta e^{\beta s}}-1)\sum_{m=1}^kY_{s-}^{(m)}\,ds\right)
\le
\exp\!\left(\thetalo(e^\eta-1)V\right),
\]
where
\[
V:=\int_{-\infty}^0 e^{\beta s}\sum_{m=1}^kY_{s-}^{(m)}\,ds.
\]
Using the stationary shot-noise representation of the sources and Tonelli's
theorem,
\[
V=\sum_{m=1}^k\int_{(-\infty,0]}(-u)e^{\beta u}\,dN_u^{(m)}
=\int_{(-\infty,0]}(-u)e^{\beta u}\,d\widetilde N_u,
\]
where $\widetilde N$ is homogeneous Poisson with rate $k\bar\mu$. The kernel
$u\mapsto(-u)e^{\beta u}\mathbf 1_{\{u\le0\}}$ is bounded and integrable, so
Campbell's formula implies $\E[e^{cV}]<\infty$ for every $c>0$. Consequently
$\E[e^{\eta W}]<\infty$ for every $\eta>0$, and in particular
$\E[W\log(1+W)]<\infty$. This proves $C_{\mathrm{init}}<\infty$ and completes
the proof.
\end{proof}

\end{document}
